\documentclass{article}

\usepackage{amsmath,amssymb}
\usepackage{xurl}
\usepackage[hidelinks]{hyperref}

% Shared commands for notes in docs/paper-gaps/.

\DeclareUnicodeCharacter{2080}{\ifmmode{}_{0}\else\textsubscript{0}\fi}
\DeclareUnicodeCharacter{2081}{\ifmmode{}_{1}\else\textsubscript{1}\fi}
\DeclareUnicodeCharacter{2082}{\ifmmode{}_{2}\else\textsubscript{2}\fi}
\DeclareUnicodeCharacter{2099}{\ifmmode{}_{n}\else\textsubscript{n}\fi}

\newcommand{\Lean}{\textsc{Lean}}
\ExplSyntaxOn
\NewDocumentCommand{\leanid}{m}{
  \ifmmode
    \textsf{\detokenize{#1}}
  \else
    \group_begin:
    \urlstyle{sf}
    \tl_set:Nn \l_tmpa_tl {#1}
    \tl_replace_all:Nnn \l_tmpa_tl {\_} {_}
    \exp_args:NV \path \l_tmpa_tl
    \group_end:
  \fi
}
\ExplSyntaxOff
\providecommand{\tr}{\operatorname{tr}}
\newcommand{\tnleanIssueBase}{https://github.com/LionSR/TNLean/issues/}
\newcommand{\tnleanPrBase}{https://github.com/LionSR/TNLean/pull/}
\newcommand{\tnleanDocBase}{https://sirui-lu.com/TNLean/docs/}
\newcommand{\ghissue}[1]{\href{\tnleanIssueBase#1}{GitHub issue \##1}}
\newcommand{\ghpr}[1]{\href{\tnleanPrBase#1}{PR \##1}}
\newcommand{\doclink}[1]{\href{\tnleanDocBase#1.html}{\path{#1}}}

% Dirac notation and the identity operator, as in blueprint/src/macros/common.tex.
% \providecommand keeps note-local definitions authoritative.
\usepackage{dsfont}
\providecommand{\ket}[1]{|#1\rangle}
\providecommand{\bra}[1]{\langle #1|}
\providecommand{\braket}[2]{\langle #1 | #2 \rangle}
\providecommand{\ketbra}[2]{|#1\rangle\!\langle #2|}
\providecommand{\Id}{\mathds{1}}
\providecommand{\one}{\mathds{1}}
\providecommand{\C}{\mathbb{C}}
\providecommand{\MN}[1]{M_{#1}(\C)}
\providecommand{\MD}{M_D(\C)}

% External referencing for paper sources.  The build helper writes sanitized
% auxiliary files under build/paper-aux/ and excludes duplicated source labels.
\usepackage{xr}

% Import a generated paper auxiliary file with a caller-supplied label prefix.
% The candidate paths support builds from docs/paper-gaps/, the repository root,
% and build/paper-aux/ itself.
\newcommand{\importpaperaux}[2]{%
  \IfFileExists{../../build/paper-aux/#2.aux}{%
    \externaldocument[#1]{../../build/paper-aux/#2}%
  }{%
    \IfFileExists{build/paper-aux/#2.aux}{%
      \externaldocument[#1]{build/paper-aux/#2}%
    }{%
      \IfFileExists{#2.aux}{%
        \externaldocument[#1]{#2}%
      }{}%
    }%
  }%
}

% General external reference macros
\newcommand{\paperref}[2]{\ref{#1-#2}}
\newcommand{\papereqref}[2]{(\ref{#1-#2})}

% CPSV16 source (1606.00608 / MPDO-22-12-17-2.tex) external document and shortcuts
\importpaperaux{cpsv16-}{1606.00608/MPDO-22-12-17-2-tnlean}
\newcommand{\cpsvref}[1]{\paperref{cpsv16}{#1}}
\newcommand{\cpsveqref}[1]{\papereqref{cpsv16}{#1}}

% Verdict marker: \gapnote{<kind>}{<status>}. Kind is the policy
% classification (clarification, local-correction, scope-restriction,
% unfaithful, false-source, open-gap); status is open, wip (elimination
% underway), resolved, or historical. Machine-read by the site index;
% prints a verdict line.
\newcommand{\gapnote}[2]{\par\noindent\textbf{Verdict:} #1 (#2).\par}


\title{The Trace-Ratio Exponent in the Blocked Vertical-Operator Formula}
\author{}
\date{2026-07-21}

\begin{document}
\maketitle
\gapnote{local-correction}{resolved}

\section{The two formulas in Appendix C.4}

Appendix~C.4 of Ref.~\cite{Cirac2016MPDO_arXiv} compares the one-site and
two-site vertical canonical forms. Lines~2001--2008 pair their sectors and
show that the corresponding representative tensors satisfy
\begin{align}
    A^{(2)}_{\sigma(\alpha)}
        &=
    \frac{m_\alpha}{n_{\sigma(\alpha)}}
    V_\alpha A^{(1)}_\alpha V_\alpha^*.
    \label{eq:cpsv16_blocked_operator_trace_ratio_exponent_letter_scaling}
\end{align}
Here
\begin{align}
    m_\alpha
        &=
    \tr(\mu_\alpha),
    \label{eq:cpsv16_blocked_operator_trace_ratio_exponent_first_multiplicity_trace}\\
    n_\beta
        &=
    \tr(\nu_\beta),
    \label{eq:cpsv16_blocked_operator_trace_ratio_exponent_second_multiplicity_trace}
\end{align}
for the positive diagonal multiplicity matrices $\mu_\alpha$ and $\nu_\beta$.
The identification of equal bond dimensions in
Eq.~\ref{eq:cpsv16_blocked_operator_trace_ratio_exponent_letter_scaling} is
harmless.

For a closed chain of length $L$, multiplying every tensor letter by a scalar
$c$ multiplies the resulting operator by $c^L$. Unitary conjugation and the
identification of bond indices leave the closing trace unchanged. Therefore
Eq.~\ref{eq:cpsv16_blocked_operator_trace_ratio_exponent_letter_scaling}
produces the coefficient
\begin{align}
    \left(\frac{m_\alpha}{n_{\sigma(\alpha)}}\right)^L
    \tr\!\left(\nu_{\sigma(\alpha)}^L\right).
    \label{eq:cpsv16_blocked_operator_trace_ratio_exponent_correct_coefficient}
\end{align}

Line~2013 of Ref.~\cite{Cirac2016MPDO_arXiv} instead prints
\begin{align}
    \frac{m_\alpha^L}{n_{\sigma(\alpha)}}
    \tr\!\left(\nu_{\sigma(\alpha)}^L\right).
    \label{eq:cpsv16_blocked_operator_trace_ratio_exponent_printed_coefficient}
\end{align}
The two expressions differ by the factor
$n_{\sigma(\alpha)}^{L-1}$. The source gives no normalization that sets this
factor to one.

\section{Confirmation later in the source}

Line~2040 of Ref.~\cite{Cirac2016MPDO_arXiv} resolves the exponent:
\begin{align}
    \tr\!\left[\left(\frac{m_\alpha}{n_{\sigma(\alpha)}}
            \nu_{\sigma(\alpha)}\right)^L\right]
        &=
    \left(\frac{m_\alpha}{n_{\sigma(\alpha)}}\right)^L
    \tr\!\left(\nu_{\sigma(\alpha)}^L\right).
    \label{eq:cpsv16_blocked_operator_trace_ratio_exponent_source_confirmation}
\end{align}
This agrees with
Eq.~\ref{eq:cpsv16_blocked_operator_trace_ratio_exponent_correct_coefficient}
and with the tensor-letter scaling at line~2008. Thus line~2013 omits the
exponent $L$ from the denominator. This is a local typographical error, not an
additional hypothesis or a change to the argument.

\section{Formal correction}

The simultaneous blocked-operator representations use
Eq.~\ref{eq:cpsv16_blocked_operator_trace_ratio_exponent_correct_coefficient}.
The formal statement retains the sector equivalence, bond-dimension
identifications, and unitary matrices that establish
Eq.~\ref{eq:cpsv16_blocked_operator_trace_ratio_exponent_letter_scaling}. It
does not normalize the multiplicity matrices.\footnote{Formal declaration:
\leanid{MPOTensor.transportedVerticalSector_exists_blockedOperatorRepresentations}
in \doclink{TNLean/MPS/MPDO/VerticalBlockedOperatorRepresentations}.}
This correction is tracked in \ghissue{4586}.

\section{Conclusion}

The coefficient consistent with the preceding tensor identity and the later
trace-power calculation is
\begin{align}
    \left(\frac{m_\alpha}{n_{\sigma(\alpha)}}\right)^L
    \tr\!\left(\nu_{\sigma(\alpha)}^L\right).
    \label{eq:cpsv16_blocked_operator_trace_ratio_exponent_verdict_coefficient}
\end{align}
The expression printed at line~2013 of
Ref.~\cite{Cirac2016MPDO_arXiv} is missing the exponent $L$ on the
denominator.

\bibliographystyle{alpha}
\bibliography{references}

\end{document}
